\section{Inferable Horizons}
\label{sec:inferable_horizons}

The horizon law is structural, but its operational use requires more than the
existence of an optimizer. A data-driven memory rule needs estimators of the
drift-geometry parameters $(H,\zeta)$, together with a quantitative criterion
that compares inferential resolution to the curvature of the useful-memory
region. This section records that bridge.

Write $\alpha = \log \zeta$, and let $m$ denote the per-lag sample size used to
estimate lagged transport discrepancies. For a fixed lag design
$1,\dots,L$, define the lag-energy information scale
\[
  \mathcal I_{m,L}(H) \coloneqq m\sum_{j=1}^L j^{2H}.
\]

\begin{assumption}[Lag-geometry inference input]
\label{ass:lag_geometry_input}
For the lag model $D_j = \zeta j^H$, suppose there are estimators
$(\widehat \alpha_{m,L},\widehat H_{m,L})$ such that, on compact parameter sets,
\[
  \sqrt{\mathcal I_{m,L}(H)}
  \begin{pmatrix}
    \widehat \alpha_{m,L} - \alpha \\
    \widehat H_{m,L} - H
  \end{pmatrix}
  \Rightarrow \mathcal N(\mathbf 0, \Sigma_{L,H,\zeta}).
\]
\end{assumption}

The power-law lag-geometry model of \citet{Parreira2026Jul} provides one
concrete source of this assumption.

The continuous optimizer of the horizon envelope can be written as the smooth
map
\[
  n_\star(\alpha,H)
  = \left(\frac{aC_K}{H C_S e^{\alpha}}\right)^{1/(a+H)},
\]
\[
  g(\alpha,H) \coloneqq \log n_\star(\alpha,H).
\]
Its gradient is
\[
  \nabla g(\alpha,H)
  =
  \begin{pmatrix}
    -(a+H)^{-1} \\
    -\bigl[H(a+H)\bigr]^{-1} - g(\alpha,H)(a+H)^{-1}
  \end{pmatrix}.
\]

\begin{theorem}[Plug-in Horizon Central Limit Law]
\label{thm:plugin_horizon_clt}
Under \Cref{ass:lag_geometry_input},
\[
  \sqrt{\mathcal I_{m,L}(H)}
  \bigl(\log \widehat n_\star - \log n_\star\bigr)
  \Rightarrow \mathcal N(0,\tau^2),
\]
where $\widehat n_\star = n_\star(\widehat \alpha_{m,L},\widehat H_{m,L})$ and
\[
  \tau^2
  = \nabla g(\alpha,H)^\top \Sigma_{L,H,\zeta} \nabla g(\alpha,H).
\]
\end{theorem}

\begin{proof}
The map $g$ is smooth on any compact set bounded away from $H=0$ and
$\zeta=0$. Apply the multivariate delta method to
$(\widehat \alpha_{m,L},\widehat H_{m,L})$ under
\Cref{ass:lag_geometry_input}.
\end{proof}

The useful-memory region is the set
\[
  \mathcal U_\delta = \{n : \Phi(n) \le (1+\delta)\Phi(n_\star)\}.
\]
Let $(x_-(\delta;a,H),x_+(\delta;a,H))$ be the normalized endpoints from
\Cref{prop:useful_memory_region}, and define the log-radius
\[
  r_\delta(a,H)
  \coloneqq
  \min\bigl\{-\log x_-(\delta;a,H),\,\log x_+(\delta;a,H)\bigr\}.
\]
This quantity measures the operational tolerance of the horizon law: it is the
largest symmetric log-error around $n_\star$ that stays inside the useful-memory
region.

\begin{theorem}[Operational Identifiability of Useful Memory]
\label{thm:operational_identifiability}
Under \Cref{ass:lag_geometry_input},
\[
  \mathbb{P}\bigl(\widehat n_\star \in \mathcal U_\delta\bigr)
  =
  2\Phi_{\mathcal N}\!\left(
    \frac{\sqrt{\mathcal I_{m,L}(H)}\,r_\delta(a,H)}{\tau}
  \right) - 1 + o(1),
\]
where $\Phi_{\mathcal N}$ is the standard normal CDF. In particular,
\[
  \frac{\sqrt{\mathcal I_{m,L}(H)}\,r_\delta(a,H)}{\tau(H,\zeta)} \to \infty
  \quad \Longrightarrow \quad
  \mathbb{P}\bigl(\widehat n_\star \in \mathcal U_\delta\bigr) \to 1.
\]
\end{theorem}

\begin{proof}
By \Cref{prop:useful_memory_region},
$\widehat n_\star \in \mathcal U_\delta$ if and only if
$|\log \widehat n_\star - \log n_\star| \le r_\delta(a,H)$. Standardize this
event with \Cref{thm:plugin_horizon_clt} and use the symmetry of the Gaussian
limit.
\end{proof}

The quantity
\[
  \mathfrak S_{\delta,m,L}
  \coloneqq
  \frac{\sqrt{\mathcal I_{m,L}(H)}\,r_\delta(a,H)}{\tau(H,\zeta)}
\]
is therefore the operational identifiability score of the horizon. It compares
lag-geometry information to the curvature-imposed precision requirement of the
useful-memory region.

Appendix-style saturation is not informative here: the relevant comparison lies
across the transition from low to moderate identifiability. Table
\ref{tab:inferable_horizon} records representative regimes across that
transition.

\begin{table}[!htbp]
\centering
\small
\caption{Representative regimes across the inferable-horizon transition.}
\label{tab:inferable_horizon}
\begin{tabular}{rrrrrrr}
\toprule
$L$ & $m$ & $H$ & $\delta$ & $\mathfrak S$ & Empirical & Gaussian \\
\midrule
5 & 20 & 0.3 & 0.050 & 0.501 & 0.367 & 0.384 \\
8 & 10 & 0.9 & 0.020 & 1.001 & 0.664 & 0.683 \\
3 & 200 & 0.3 & 0.010 & 1.996 & 0.949 & 0.954 \\
8 & 100 & 0.3 & 0.005 & 3.001 & 0.997 & 0.997 \\
\bottomrule
\end{tabular}
\end{table}


\begin{proposition}[Plug-in Horizon Regret]
\label{prop:plugin_horizon_regret}
Let $u = \log(\widehat n_\star/n_\star)$. Then, as $u \to 0$,
\[
  \frac{\Phi(\widehat n_\star)-\Phi(n_\star)}{\Phi(n_\star)}
  = \frac{aH}{2}u^2 + O(u^3).
\]
Consequently,
\[
  \mathbb E\!
  \left[
    \frac{\Phi(\widehat n_\star)-\Phi(n_\star)}{\Phi(n_\star)}
  \right]
  = \frac{aH}{2}\frac{\tau^2}{\mathcal I_{m,L}(H)} + o\!\left(\mathcal I_{m,L}(H)^{-1}\right).
\]
\end{proposition}

\begin{proof}
By \Cref{prop:useful_memory_region},
\[
  \frac{\Phi(n_\star e^u)}{\Phi(n_\star)} = \Psi(e^u)
  = \frac{H e^{-au} + a e^{Hu}}{a+H}.
\]
Differentiate at $u=0$. The first derivative vanishes and the second derivative
equals $aH$, giving the quadratic expansion. The expectation bound then follows
from \Cref{thm:plugin_horizon_clt} and the boundedness of third moments on
compact parameter sets.
\end{proof}

\begin{figure}[!htbp]
\centering
\safeincludegraphics[width=0.94\linewidth]{artifacts/figures/inferable_horizon/fig_operational_identifiability}
\caption{Operational identifiability of useful memory. Across the tested grid,
useful-region hit rates collapse against the score
$\mathfrak S_{\delta,m,L} = \sqrt{\mathcal I_{m,L}(H)}r_\delta(a,H)/\tau(H,\zeta)$,
matching the Gaussian approximation from \Cref{thm:operational_identifiability}.}
\label{fig:inferable_horizon}
\end{figure}

This section separates two questions that need not agree. A temporal-validity
horizon may exist structurally without being operationally identifiable from the
available lag geometry. The gap between existence and operational inferability
is governed by the score $\mathfrak S_{\delta,m,L}$.
